\chapter{Il DSL generico: un linguaggio per i vincoli}
\label{ch:metodo_dsl}

\epigraph{``A domain-specific language is a computer
language specialized to a particular application
domain.''}{Martin Fowler, \emph{Domain-Specific
Languages}~\parencite{fowler2010}}

\lettrine[lines=3]{P}{rima} o poi tutti i sistemi di
scheduling devono affrontare lo stesso problema editoriale:
gli utenti vogliono esprimere vincoli che non sono mai
stati previsti. ``Mai due ore di matematica consecutive nei
gioved\`i pari delle classi prime''; ``almeno un'ora di
laboratorio per ogni classe di scientifico''; ``il prof
Bianchi e la prof Verdi mai stesso slot stessa giornata''.
\`E tecnicamente impossibile prevedere tutto e codificarlo
hard-coded nel solver.

La soluzione classica \`e un \emph{Domain-Specific
Language}\index{DSL}: un piccolo linguaggio dichiarativo che
gli utenti possono usare per scrivere i propri vincoli. In
$\pi$Tantum il DSL si chiama \emph{DSL generico} e lo
descriviamo in questo capitolo dal punto di vista
implementativo.

\section{Filosofia: ``un parser, molti compilatori''}

\begin{ideabox}
Il DSL generico ha una grammatica unica e un parser unico,
ma ha pi\`u \emph{back-end}: lo stesso vincolo si pu\`o
\emph{valutare} sui dati correnti (per check sincroni o per
filtri lista), si pu\`o \emph{compilare} a vincoli CP-SAT
(per Phase~B), oppure si pu\`o \emph{compilare} a delta SOFT
per il drag-and-drop in tempo reale dell'editor Orario.
\`E lo schema classico ``parser $\to$ AST $\to$ molti
compilatori'' del Dragon Book~\parencite{aho2006}.
\end{ideabox}

\section{Grammatica e parser (rinvio)}

La grammatica BNF completa --- con le otto sorgenti
iterabili (\texttt{teachers}, \texttt{classes},
\texttt{subjects}, \texttt{classrooms}, \texttt{groups},
\texttt{students}, \texttt{lessons}, \texttt{slots}) e i path
arbitrari del tipo \texttt{l.classroom.tags} --- vive nel
capitolo tutorial alla
sezione~\ref{sec:dsl-bnf}. Qui basti dire che il
parser \`e \emph{recursive descent} a discesa ricorsiva, scritto
in Python puro senza generatori (no PLY/Lark) e che le
precedenze vanno da \texttt{<=>} (la pi\`u bassa) fino al
\texttt{.} di accesso a membro (la pi\`u alta), passando per
\texttt{=>}, \texttt{or}, \texttt{and}, \texttt{not} e gli
operatori di confronto.

\section{AST}

\begin{defbox}
L'AST (\emph{Abstract Syntax Tree})\index{AST} prodotto dal
parser ha nodi che sono piccole dataclass:
\texttt{Quant}, \texttt{Count}, \texttt{BinOp},
\texttt{Cmp}, \texttt{Ref}, \texttt{Lit},
\texttt{Builtin}. Riconosce token tramite
\texttt{tokenize()} con regex, e poi assembla i nodi per
discesa ricorsiva.
\end{defbox}

Esempio: l'espressione
\begin{lstlisting}
forall t in teachers where t.subject == "Fisica":
  count l in lessons where l.teacher == t and
                            l.classroom.type == "lab_fisica":
    l == 1
\end{lstlisting}

produce questo AST (semplificato):
\begin{lstlisting}
Quant(kind='forall', var='t', source='teachers',
      filter=Cmp('==', Ref(['t', 'subject']), Lit('Fisica')),
      body=Count(var='l', source='lessons',
                 filter=And([
                   Cmp('==', Ref(['l', 'teacher']), Ref(['t'])),
                   Cmp('==', Ref(['l', 'classroom', 'type']),
                            Lit('lab_fisica'))]),
                 cmp=Cmp('==', Var('count'), Lit(1))))
\end{lstlisting}

\section{Tre back-end}

\subsection{Evaluator: ``vale o no?''}

\begin{defbox}
L'\emph{evaluator}\index{evaluator (DSL)} prende un AST e
i dati correnti del sistema (\texttt{teachers},
\texttt{lessons}, ecc.) e restituisce \texttt{True/False}
per i vincoli \texttt{forall}/\texttt{exists}, e un intero
per i \texttt{count}.
\end{defbox}

\`E usato in modalit\`a sincrona per i filtri lista (``mostrami
tutti i docenti per cui questo vincolo \`e violato''), per i
check pre-pipeline (verifica di consistenza), e come oracolo
nei test.

\subsection{Compilatore CP-SAT}

\begin{defbox}
Il \emph{compilatore CP-SAT} traduce l'AST in vincoli OR-Tools.
Le quantificazioni diventano cicli sui domini; le
comparazioni diventano vincoli aritmetici; i \texttt{count}
diventano somme con vincoli di uguaglianza/disuguaglianza.
\end{defbox}

Esempio: \texttt{count l in lessons where l.teacher == t: l == 1}
diventa
\begin{lstlisting}[style=py]
total = sum(BoolVar(f"l_{l.id}_active") for l in lessons
            if l.teacher_id == t.id)
model.Add(total == 1)
\end{lstlisting}

\subsection{Compilatore delta SOFT}

\begin{defbox}
Il \emph{compilatore delta} produce, dato un vincolo SOFT e
una proposta di mossa (drag-and-drop di una lezione), il
\emph{delta} dell'obiettivo (positivo se la mossa peggiora,
negativo se migliora) in $O(1)$.
\end{defbox}

\`E quello che permette all'editor di mostrare in tempo reale
``+0.3 SOFT'' o ``$-0.7$ SOFT'' mentre l'utente trascina una
lezione.

\section{Estensione: aggiungere un built-in}

Aggiungere un nuovo built-in (es.\ \texttt{distinct\_days}) \`e
un'operazione locale al modulo
\texttt{webui/backend/utils/general\_dsl.py} e tocca quattro punti:
\texttt{tokenize()}, parser, evaluator e compilatore CP-SAT
(quest'ultimo di norma 5--10 righe). La
sezione~\ref{sec:dsl-estendere} del capitolo tutorial sviluppa
la procedura passo-passo --- con i casi simmetrici
\emph{nuova sorgente iterabile}, \emph{nuovo pragma} e
\emph{nuovo operatore} --- e mostra il test di accettazione in
\texttt{test\_general\_dsl.py}.

\section{Performance}

\begin{ideabox}
Il parser \`e veloce ($\sim 10\mu$s per espressione tipica
da 50 token). L'evaluator \`e proporzionale al prodotto dei
domini (es. $|\text{teachers}| \cdot |\text{lessons}|$ per il
vincolo dell'esempio). Il compilatore CP-SAT genera
tipicamente $\sim 100$ vincoli per ogni vincolo DSL non
banale; per scuole grandi, l'aggiunta di 20--30 vincoli DSL
non rallenta significativamente la pipeline.
\end{ideabox}

\section{Roadmap}

Funzionalit\`a non ancora implementate ma pianificate:
\begin{itemize}
\item \textbf{Live validation}: squiggle rosso sotto la
  parte invalida, suggerimenti fuzzy, quick-fix.
\item \textbf{Autocomplete}: nel campo DSL, lista delle
  sorgenti / attributi / built-in con descrizione.
\item \textbf{Hover docs}: passando il mouse su un
  identifier, mostra documentazione contestuale.
\item \textbf{Natural-language preview}: traduzione
  automatica del vincolo DSL in italiano leggibile.
\end{itemize}

\section{Connessioni}

\begin{connessionebox}
\begin{itemize}
\item \textbf{CP-SAT} (cap.~\ref{ch:metodo_cpsat}): il
  back-end principale per i vincoli HARD.
\item \textbf{Editor Orario}: il compilatore delta abilita
  il drag-and-drop con preview live.
\item \textbf{Vincoli logici DNF}
  (cap.~\ref{ch:dsl_generico_per_i_vincoli}): il DSL
  generico \`e una generalizzazione del vecchio sistema di
  ``vincoli logici disgiuntivi'' che ne preserva la
  semantica come caso particolare.
\end{itemize}
\end{connessionebox}

\begin{biblioparvabox}
\begin{itemize}
\item \cite{aho2006}: il \emph{Dragon Book}, riferimento
  classico per parser, AST, compilatori.
\item \cite{fowler2010}: \emph{Domain-Specific Languages},
  approccio pratico ai DSL.
\item \cite{marriott1998}: \emph{Programming with
  Constraints}, fondamenti della programmazione a
  vincoli.
\end{itemize}
\end{biblioparvabox}
